\documentclass[11pt,reqno]{amsart}
\usepackage{amsmath,amssymb,amsthm}
\usepackage{graphicx}
\usepackage{booktabs}
\usepackage{geometry}
\usepackage{subcaption}
\usepackage{hyperref}
\usepackage{microtype}
\geometry{margin=1in}

\newtheorem{theorem}{Theorem}
\newtheorem{conjecture}[theorem]{Conjecture}
\newtheorem{observation}[theorem]{Observation}
\newtheorem{definition}[theorem]{Definition}
\newtheorem{remark}[theorem]{Remark}

\title{Finite-$n$ Convergence and Period-4 Oscillatory Structure\\
in Hook Statistics of Self-Conjugate Partitions}

\author{Rudra Mishra}
\address{Independent Researcher}
\email{}
\date{June 6, 2026}

\begin{document}

\begin{abstract}
We present an experimental investigation of the finite-$n$ 
convergence structure of the Craig--Ono--Singh asymptotic formula 
for $t$-hook statistics in self-conjugate integer partitions. 
Craig, Ono, and Singh proved that for fixed $t \geq 1$, the number 
of $t$-hooks among self-conjugate partitions of $n$ is 
asymptotically normally distributed with explicit mean 
$\mu_t(n) \sim \frac{\sqrt{6n}}{\pi} - \frac{t}{2} + 
\frac{3}{\pi^2} + \frac{\delta_t}{4}$. 
While their work establishes large-$n$ asymptotics, the 
finite-$n$ error structure has not been computationally 
characterized. For $t = 1, 2, 3$ and $3 \leq n \leq 60$, 
we define the signed error $e(t,n) = \mu_t^{\mathrm{emp}}(n) - 
\mu_t^{\mathrm{theory}}(n)$. The scaled error $e(t,n)\sqrt{n}$ 
remains bounded, numerically verifying the theoretical $O(n^{-1/2})$ bound 
for the tested range. Furthermore, exploratory Fast Fourier transform analysis 
detects a frequency near $0.25$ (period-4) in the SC error spectrum. This aligns 
with the expected analytic contributions from subdominant singularities at fourth 
roots of unity in the generating function. We propose a descriptive empirical 
ansatz for this period-4 oscillatory structure and outline the necessary statistical 
and analytic directions required to formalize it for larger $n$.
\end{abstract}

\maketitle

\tableofcontents

%---------------------------------------------------------------
\section{Introduction}
%---------------------------------------------------------------

The study of integer partitions sits at the crossroads of 
combinatorics, number theory, and mathematical physics. 
A partition of a positive integer $n$ is a nonincreasing 
sequence of positive integers summing to $n$. Among the 
various combinatorial statistics attached to partitions, 
hook lengths occupy a distinguished role: they govern the 
representation theory of symmetric groups via the 
Frame--Thrall--Robinson hook length formula, appear in the 
Nekrasov--Okounkov identity, and connect to modular forms 
through the work of Han.

A partition $\lambda$ is \emph{self-conjugate} if its Young 
diagram equals its transpose. Self-conjugate partitions are 
in natural bijection with partitions into distinct odd parts 
via hook decomposition, a classical fact going back to Euler. 
Diagonal hooks in self-conjugate partitions are odd, though general 
hook lengths may be even or odd. This creates a structural 
constraint that distinguishes them sharply from unrestricted partitions.

Craig, Ono, and Singh~\cite{COS} recently established that 
the number of $t$-hooks in self-conjugate partitions of $n$ 
is asymptotically normally distributed, extending earlier 
work of Griffin, Ono, and Tsai on unrestricted partitions. 
Their proof uses two-variable generating functions derived 
via the Littlewood bijection for $t$-core partitions, 
together with Euler--Maclaurin summation, the dilogarithm 
function, and the saddle-point method.

\medskip

The present work operates in a different regime. Rather than 
proving new asymptotic theorems, we conduct an exploratory computational investigation 
of the finite-$n$ remainder term of the Craig--Ono--Singh formula. 
Specifically, we ask: how does the empirical mean $\mu_t^{\mathrm{emp}}(n)$ approach 
$\mu_t^{\mathrm{theory}}(n)$ as $n$ grows?

The Craig--Ono--Singh asymptotic contains an uncharacterized remainder 
term of order $O(n^{-1/2})$. Consequently, the empirical residual studied 
here mixes finite-$n$ fluctuations with the unknown higher-order 
asymptotic correction. 

Our primary empirical observation is the numerical detection of a period-4 
oscillatory component in the residual for the tested range $n \le 60$. Rather 
than an unexpected anomaly, this structure is the anticipated analytic consequence 
of the arithmetic of self-conjugate Young diagrams, specifically arising from 
singularities at fourth roots of unity in the generating function.

\medskip

\noindent\textbf{Organization.} Section~2 provides background 
on partitions, Young diagrams, hooks, and the Craig--Ono--Singh 
theorem. Section~3 describes our computational framework. 
Section~4 presents experimental observations: scaled error, 
FFT analysis, mod-4 decomposition, and model fitting. 
Section~5 compares SC and unrestricted partitions. 
Section~6 states the empirical model. Section~7 discusses 
methodological limitations and future directions.

%---------------------------------------------------------------
\section{Background}
%---------------------------------------------------------------

\subsection{Partitions and Young diagrams}

A \emph{partition} of $n$ is a sequence 
$\lambda = (\lambda_1 \geq \lambda_2 \geq \cdots \geq \lambda_k > 0)$ 
with $\sum_i \lambda_i = n$. We write $\lambda \vdash n$ and 
denote the number of partitions of $n$ by $p(n)$.

The \emph{Young diagram} of $\lambda$ is the left-justified 
array of boxes with $\lambda_i$ boxes in row $i$. The 
\emph{conjugate partition} $\lambda'$ is obtained by 
transposing this array.

\subsection{Self-conjugate partitions}

A partition is \emph{self-conjugate} if $\lambda = \lambda'$, 
i.e., the Young diagram is symmetric across its main diagonal. 
We denote by $\mathrm{SC}(n)$ the set of self-conjugate 
partitions of $n$ and by $\mathrm{sc}(n) = |\mathrm{SC}(n)|$ 
their count.

\begin{observation}[Euler, 1748]
$\mathrm{sc}(n)$ equals the number of partitions of $n$ into 
distinct odd parts.
\end{observation}

This is proved via hook decomposition: the diagonal hooks of 
a self-conjugate partition have odd lengths, and stripping them 
yields a bijection with distinct odd part partitions.

\subsection{Hook lengths}

For a box $(i,j)$ in the Young diagram of $\lambda$, the 
\emph{hook length} is
\[
h(i,j) = (\lambda_i - j) + (\lambda'_j - i) + 1.
\]
The \emph{hook multiset} $H(\lambda)$ collects all hook 
lengths in $\lambda$.

For a fixed integer $t \geq 1$, define
\[
N_t(\lambda) = \#\{h \in H(\lambda) : h = t\}.
\]

In self-conjugate partitions, the diagonal hooks are odd; however, 
general hook lengths can be even or odd.

\subsection{The Craig--Ono--Singh theorem}

The following is the main result of~\cite{COS}.

\begin{theorem}[Craig--Ono--Singh, 2025]\label{thm:COS}
Let $t \geq 1$ and let $\widehat{N}_{t,n}$ be the random 
variable giving the distribution of $N_t(\lambda)$ over 
$\mathrm{SC}(n)$. Then as $n \to \infty$, 
$\widehat{N}_{t,n}$ is asymptotically normally distributed 
with mean
\[
\mu_t(n) = \frac{\sqrt{6n}}{\pi} - \frac{t}{2} 
+ \frac{3}{\pi^2} + \frac{\delta_t}{4} 
+ O\!\left(\frac{1}{\sqrt{n}}\right)
\]
and variance
\[
\sigma_t(n)^2 = \frac{(\pi^2 - 6)\sqrt{6n}}{\pi^3} + O(1),
\]
where $\delta_t = 1$ if $t$ is odd and $0$ otherwise.
\end{theorem}

%---------------------------------------------------------------
\section{Computational Framework}
%---------------------------------------------------------------

\subsection{Implementation}

All computations were performed in Python on standard hardware. 
The codebase includes:
\begin{itemize}
\item recursive partition generation,
\item self-conjugacy testing via conjugate comparison,
\item hook length computation for all boxes,
\item empirical mean calculation per $n$ and $t$,
\item FFT analysis via \texttt{numpy.fft},
\item least-squares model fitting via \texttt{scipy.optimize.curve\_fit},
\item entropy calculation for hook distributions.
\end{itemize}

\subsection{Definitions}

Fix $t \geq 1$. For each $n$, define the \emph{empirical mean}
\[
\mu_t^{\mathrm{emp}}(n) 
= \frac{1}{\mathrm{sc}(n)} 
\sum_{\lambda \in \mathrm{SC}(n)} N_t(\lambda).
\]

Define the \emph{signed error}
\[
e(t,n) = \mu_t^{\mathrm{emp}}(n) - \mu_t^{\mathrm{theory}}(n),
\]
where $\mu_t^{\mathrm{theory}}(n)$ is the leading asymptotic 
from Theorem~\ref{thm:COS}.

For \emph{unrestricted} partitions, the theoretical mean from Griffin--Ono--Tsai~\cite{GOT} is
\[
\mu_t^{\mathrm{GOT}}(n) = \frac{\sqrt{6n}}{\pi} - \frac{t}{2} + \frac{3}{\pi^2} + O\!\left(\frac{1}{\sqrt{n}}\right).
\]
Note the absence of the $\delta_t/4$ term compared to the self-conjugate case. We define the signed error for unrestricted partitions analogously using $\mu_t^{\mathrm{GOT}}(n)$.

Define the \emph{relative error}
\[
R(t,n) = \frac{|e(t,n)|}{|\mu_t^{\mathrm{theory}}(n)|}.
\]

Define the \emph{scaled error}
\[
A(t,n) = e(t,n)\sqrt{n}.
\]

\subsection{Computational range}

Experiments were run for $t \in \{1, 2, 3\}$ and 
$3 \leq n \leq 60$. For comparison, unrestricted partition 
hook statistics were computed over the same range.

Figure~\ref{fig:master} illustrates the growth of $p(n)$, 
$p_{\mathrm{odd}}(n)$, $p_{\mathrm{distinct}}(n)$, and 
$\mathrm{sc}(n)$ for $n \leq 20$, empirically confirming Euler's theorem 
($p_{\mathrm{odd}} = p_{\mathrm{distinct}}$) and the 
self-conjugate--distinct odd parts bijection.

\begin{figure}[htbp]
\centering
\includegraphics[width=0.85\textwidth]{figure1_partition_sequences.png}
\caption{Growth of partition sequences. Red and green 
curves (p\_odd and p\_distinct) are visually identical, 
empirically confirming Euler's theorem.}
\label{fig:master}
\end{figure}

%---------------------------------------------------------------
\section{Experimental Observations}
%---------------------------------------------------------------

\subsection{Euler's theorem and Young diagram bijection}

As a verification baseline, Figure~\ref{fig:euler} indicates 
that empirical mean hook length is consistently lower for 
self-conjugate partitions than for all partitions of $n$, 
consistent with the geometric compactness of self-conjugate 
Young diagrams.

\begin{figure}[htbp]
\centering
\includegraphics[width=0.75\textwidth]{figure2_hook_stats.png}
\caption{Mean hook length for all partitions vs.\ 
self-conjugate partitions. The gap is consistently 
negative, reflecting the geometric compactness of 
self-conjugate diagrams.}
\label{fig:euler}
\end{figure}

\subsection{Empirical vs.\ asymptotic mean}

Figure~\ref{fig:asymp} illustrates the empirical mean 
$\mu_t^{\mathrm{emp}}(n)$ alongside the Craig--Ono--Singh 
asymptotic $\mu_t^{\mathrm{theory}}(n)$ for $t = 1, 2, 3$.

\begin{figure}[htbp]
\centering
\includegraphics[width=\textwidth]{figureA_empirical_vs_asymptotic.png}
\caption{Empirical mean (solid) vs.\ Craig--Ono--Singh 
asymptotic (dashed) for $t=1,2,3$. Convergence is visible 
but oscillatory.}
\label{fig:asymp}
\end{figure}

Figure~\ref{fig:diff} indicates the signed difference 
$e(t,n)$ for all three values of $t$.

\begin{figure}[htbp]
\centering
\includegraphics[width=0.85\textwidth]{figureB_convergence_diff.png}
\caption{Signed error $e(t,n)$ for $t=1,2,3$. 
The error oscillates around zero rather than 
decaying monotonically.}
\label{fig:diff}
\end{figure}

\subsection{Relative error decay}

Figure~\ref{fig:relErr} illustrates $R(t,n)$ for $t=1,2,3$. 
Even $t$ exhibits a different decay profile than odd $t$. 
Because the $\delta_t/4$ constant is already subtracted from the 
theoretical mean, this difference confirms that the remaining 
higher-order asymptotic terms also depend strongly on the parity of $t$.

\begin{figure}[htbp]
\centering
\includegraphics[width=0.85\textwidth]{lens1_relative_error.png}
\caption{Relative error $R(t,n)$ and parity-of-$n$ split. 
Odd $t$ exhibits persistently higher relative error for 
$n < 30$.}
\label{fig:relErr}
\end{figure}

\subsection{Scaled error and decay rate}

Figure~\ref{fig:scaled} indicates the scaled error 
$A(t,n) = e(t,n)\sqrt{n}$. The quantity remains bounded 
in the tested range, rather than decaying to zero or growing. This 
does not constitute a new bound, but successfully verifies the 
$O(n^{-1/2})$ decay rate established in Theorem~\ref{thm:COS} 
for the finite regime $n \le 60$.

\begin{figure}[htbp]
\centering
\includegraphics[width=0.85\textwidth]{lens5_scaled_error.png}
\caption{Scaled error $e(t,n)\sqrt{n}$ for $t=1,2,3$. 
The quantity is bounded and oscillatory, numerically verifying the 
theoretical $O(n^{-1/2})$ decay.}
\label{fig:scaled}
\end{figure}

Figure~\ref{fig:decay} indicates log $R(t,n)$. The 
roughly linear downward trend suggests polynomial 
rather than exponential decay.

\begin{figure}[htbp]
\centering
\includegraphics[width=0.75\textwidth]{lens4_decay_rate.png}
\caption{Log relative error. The downward linear trend 
suggests polynomial decay, not exponential.}
\label{fig:decay}
\end{figure}

\subsection{FFT analysis and period-4 structure}

Figure~\ref{fig:fft} illustrates the raw FFT of $e(t,n)$. 
A dominant peak appears near frequency $0.25$, 
suggesting a period-4 oscillation. This peak 
appears to persist after parity averaging (Figure~\ref{fig:fftdecomp}), 
ruling out pure period-2 (parity) as the sole source. 

It must be noted that the raw signed error contains a non-stationary 
decaying trend. Applying an FFT to a short sequence of length 58 without 
detrending or tapering risks spectral leakage. The peak at 0.25 is therefore 
an exploratory observation; future rigorous analysis on larger datasets must employ 
proper detrending to formally isolate the periodic signal.

\begin{figure}[htbp]
\centering
\includegraphics[width=0.85\textwidth]{lens2_signed_fft.png}
\caption{Left: signed error $e(t,n)$. Right: Raw FFT of 
$e(t,n)$. The dominant frequency is near $0.25$ 
(period-4).}
\label{fig:fft}
\end{figure}

\begin{figure}[htbp]
\centering
\includegraphics[width=\textwidth]{test4_fft_decomposition.png}
\caption{FFT of raw error, scaled error, and 
parity-averaged error. The period-4 peak at frequency 
$0.25$ persists after parity averaging.}
\label{fig:fftdecomp}
\end{figure}

\subsection{Mod-4 residue analysis}

Figure~\ref{fig:mod4} separates $e(t,n)$ by residue 
class of $n$ modulo 4.

\begin{figure}[htbp]
\centering
\includegraphics[width=\textwidth]{mod4_split.png}
\caption{Signed error separated by $n \bmod 4$. 
The four residue classes exhibit distinct behavior 
for SC partitions.}
\label{fig:mod4}
\end{figure}

Table~\ref{tab:mod4} indicates the mean of $A(t,n) = 
e(t,n)\sqrt{n}$ stratified by $n \bmod 4$, for 
both SC and unrestricted partitions.

\begin{table}[htbp]
\centering
\begin{tabular}{llrrrr}
\toprule
$t$ & Type & $n\equiv 0$ & $n\equiv 1$ & 
      $n\equiv 2$ & $n\equiv 3$ \\
\midrule
1 & SC           & $-0.440$ & $-0.133$ & $\phantom{-}1.296$ & $\phantom{-}0.544$ \\
1 & Unrestricted & $\phantom{-}1.810$ & $\phantom{-}1.859$ & $\phantom{-}1.826$ & $\phantom{-}1.847$ \\
\midrule
2 & SC           & $\phantom{-}1.567$ & $\phantom{-}1.065$ & $-0.775$ & $\phantom{-}0.105$ \\
2 & Unrestricted & $\phantom{-}2.161$ & $\phantom{-}2.049$ & $\phantom{-}2.168$ & $\phantom{-}2.020$ \\
\midrule
3 & SC           & $\phantom{-}0.125$ & $\phantom{-}0.516$ & $\phantom{-}1.862$ & $\phantom{-}1.269$ \\
3 & Unrestricted & $\phantom{-}2.513$ & $\phantom{-}2.566$ & $\phantom{-}2.564$ & $\phantom{-}2.593$ \\
\bottomrule
\end{tabular}
\caption{Mean of $e(t,n)\sqrt{n}$ by residue class 
$n \bmod 4$. Given the small sample size per residue class 
(approx. 15 points), these means lack standard errors and 
formal ANOVA testing; they are presented as descriptive statistics 
rather than proven mod-4 effects.}
\label{tab:mod4}
\end{table}

\subsection{Oscillatory model fit}

We fit the model
\[
e(t,n) \approx 
\frac{A_t + B_t\cos(\pi n) 
+ C_t\cos(\pi n/2) 
+ D_t\sin(\pi n/2)}{\sqrt{n}}
\]
via least squares. Table~\ref{tab:fit} presents fitted 
constants and RMS residuals. 

\paragraph{Note on model fitting.} This is a saturated, overparameterized fit for a sequence of length 58. Without cross-validation or formal F-tests comparing it to a reduced trend-only model, the constants $C_t, D_t$ risk overfitting. The model serves solely as a descriptive ansatz for $n \le 60$.

\begin{table}[htbp]
\centering
\begin{tabular}{llrrrrr}
\toprule
$t$ & Type & $A_t$ & $B_t$ & $C_t$ & $D_t$ & RMS \\
\midrule
1 & SC           & $0.390$ & $\phantom{-}0.245$ & $-1.171$ & $-0.416$ & $0.213$ \\
1 & Unrestricted & $1.402$ & $-0.030$ & $-0.067$ & $\phantom{-}0.079$ & $0.123$ \\
2 & SC           & $0.400$ & $-0.210$ & $\phantom{-}1.651$ & $\phantom{-}0.768$ & $0.197$ \\
2 & Unrestricted & $1.628$ & $\phantom{-}0.168$ & $-0.048$ & $\phantom{-}0.126$ & $0.137$ \\
3 & SC           & $0.895$ & $\phantom{-}0.175$ & $-1.075$ & $-0.556$ & $0.277$ \\
3 & Unrestricted & $2.102$ & $-0.058$ & $-0.140$ & $-0.031$ & $0.134$ \\
\bottomrule
\end{tabular}
\caption{Fitted constants for the descriptive oscillatory 
model for $n \le 60$.}
\label{tab:fit}
\end{table}

Figure~\ref{fig:fit} illustrates the model fit overlaid 
on the empirical error.

\begin{figure}[htbp]
\centering
\includegraphics[width=\textwidth]{oscillatory_fit.png}
\caption{Empirical error (solid) and fitted 
oscillatory ansatz (dashed) for $t=1,2,3$.}
\label{fig:fit}
\end{figure}

%---------------------------------------------------------------
\section{Comparison: SC vs.\ Unrestricted Partitions}
%---------------------------------------------------------------

Figures~\ref{fig:comp1}--\ref{fig:comp3} present the 
direct comparison between SC and unrestricted partitions.

\begin{figure}[htbp]
\centering
\includegraphics[width=\textwidth]{comp1_signed_error.png}
\caption{Signed error for SC (solid) vs.\ 
unrestricted (dashed) partitions.}
\label{fig:comp1}
\end{figure}

\begin{figure}[htbp]
\centering
\includegraphics[width=\textwidth]{comp2_scaled_error.png}
\caption{Scaled error $e(t,n)\sqrt{n}$ for SC 
vs.\ unrestricted. The unrestricted scaled error 
drifts upward.}
\label{fig:comp2}
\end{figure}

\begin{figure}[htbp]
\centering
\includegraphics[width=\textwidth]{comp3_fft.png}
\caption{FFT comparison. SC partitions (solid) 
exhibit a peak near frequency $0.25$. 
Unrestricted partitions (dashed) exhibit a nearly 
flat spectrum.}
\label{fig:comp3}
\end{figure}

The contrast in Table~\ref{tab:mod4} is sharp, but requires careful interpretation. 
The unrestricted asymptotic mean used here is only a leading term; the actual error 
in the unrestricted case is known to be $O(1)$, not $O(n^{-1/2})$. The observed upward drift 
of the unrestricted scaled error (Figure~\ref{fig:comp2}) is exactly this $O(1)$ 
offset being multiplied by $\sqrt{n}$. Because this dominant non-oscillatory offset masks 
finer structure, the absence of a period-4 peak in unrestricted partitions should be interpreted 
cautiously. To make the comparison fully meaningful in future work, the $O(1)$ constant 
must be extracted from the unrestricted asymptotic before raw periodicities are compared.

%---------------------------------------------------------------
\section{Empirical Model}
%---------------------------------------------------------------

The experimental evidence motivates the following empirical fit.

\begin{observation}\label{obs:main}
Let $t \geq 1$ be fixed. Least-squares fitting suggests the descriptive empirical 
ansatz for the signed error $e(t,n) = \mu_t^{\mathrm{emp}}(n) - \mu_t^{\mathrm{theory}}(n)$ 
for self-conjugate partitions in the tested range $n \le 60$:
\[
e(t,n)\sqrt{n} \approx A_t + B_t(-1)^n 
+ C_t\cos\!\left(\frac{\pi n}{2}\right) 
+ D_t\sin\!\left(\frac{\pi n}{2}\right) 
\]
for real constants $A_t, B_t, C_t, D_t$ depending only on $t$. 
\end{observation}

\begin{remark}
The period-4 terms $\cos(\pi n/2)$ and $\sin(\pi n/2)$ 
take the four values $(1,0,-1,0)$ and $(0,1,0,-1)$ on 
residue classes $n \equiv 0,1,2,3 \pmod{4}$ respectively, 
naturally producing the mod-4 amplitude structure observed 
in Table~\ref{tab:mod4}.
\end{remark}

\begin{remark}
Rather than a purely unexpected discovery, the period-4 signal is an anticipated 
analytic feature. The generating function for self-conjugate partitions involves 
products of the form $(-q;q^2)_\infty$. When expanded near fourth roots 
of unity, the subdominant singularities at $q = \pm i$ naturally give rise 
to terms involving $\cos(\pi n/2)$ and $\sin(\pi n/2)$ via the circle method. 
Our numerical data serves to verify the presence of these expected subdominant 
contributions. The empirical model assumes constant amplitudes $C_t, D_t$ in the scaled error, 
which serves as a descriptive ansatz for this finite range, though the true 
asymptotic amplitudes may exhibit slight $n$-dependence.
\end{remark}

%---------------------------------------------------------------
\section{Discussion and Future Directions}
%---------------------------------------------------------------

\subsection{Summary of findings}

\begin{enumerate}
\item The scaled error $e(t,n)\sqrt{n}$ remains bounded, 
numerically verifying the $O(n^{-1/2})$ error bound established by 
Craig--Ono--Singh for the tested range $n \le 60$.

\item Exploratory FFT analysis and mod-4 residue decomposition 
detect a period-4 oscillatory component in the finite-$n$ residual.

\item This period-4 structure aligns with the analytic contributions 
expected from subdominant singularities at fourth roots of unity in the 
generating function.

\item Even $t$ empirically observes a different decay profile 
than odd $t$, confirming that higher-order asymptotic terms 
depend strongly on parity.
\end{enumerate}

\subsection{Limitations \& Future directions}

These results are purely computational and exploratory. 
The range $n \leq 60$ is modest, and the statistical fitting risks 
overparameterization. To advance this exploratory note into rigorous 
asymptotic theory, the following steps are necessary:

\begin{itemize}
\item \textbf{Larger computational range.} Extend computations 
to $n \geq 10^3$ using optimized partition generation via the distinct 
odd parts bijection to verify if the periodic amplitudes stabilize.

\item \textbf{Rigorous Statistical Testing.} Re-evaluate the FFT using 
proper detrending and windowing to avoid non-stationary spectral leakage, 
and apply formal ANOVA testing to the residue classes.

\item \textbf{Theoretical derivation.} Extract the exact next term in the 
asymptotic expansion using the circle method near $q = \pm i$, and compare the 
predicted oscillatory coefficients directly with the empirical ones.

\item \textbf{Refined Unrestricted Comparison.} Extract the $O(1)$ constant from the 
unrestricted asymptotic to allow for an unbiased comparison of periodic residuals 
between SC and unrestricted partitions.
\end{itemize}

%---------------------------------------------------------------
\section*{Acknowledgements}
%---------------------------------------------------------------

The author thanks the computational tools Python, NumPy, 
SciPy, and Matplotlib. This work was conducted 
independently as an exploratory research project.

%---------------------------------------------------------------
\begin{thebibliography}{9}
%---------------------------------------------------------------

\bibitem{COS}
W.~Craig, K.~Ono, and A.~Singh,
\textit{Distribution of hooks in self-conjugate partitions},
Discrete Math.\ \textbf{348} (2025).
\href{https://arxiv.org/abs/2406.09059}{arXiv:2406.09059}

\bibitem{AAOS}
T.~Amdeberhan, G.~E.~Andrews, K.~Ono, and A.~Singh,
\textit{Hook lengths in self-conjugate partitions},
Proc.\ Amer.\ Math.\ Soc.\ Ser.\ B \textbf{11} (2024), 345--357.

\bibitem{GOT}
M.~Griffin, K.~Ono, and W.-L.~Tsai,
\textit{Distributions of hook lengths in integer partitions},
Proc.\ Amer.\ Math.\ Soc., to appear.

\bibitem{Euler}
L.~Euler,
\textit{Introductio in analysin infinitorum},
1748.

\end{thebibliography}

\end{document}
